\documentclass[11pt]{article}

\usepackage[margin=1in]{geometry}
\usepackage{amsmath,amssymb,amsthm}
\usepackage{graphicx}
\usepackage{hyperref}

\newtheorem{theorem}{Theorem}
\newtheorem{lemma}{Lemma}
\newtheorem{remark}{Remark}

\newcommand{\F}{\mathcal F}
\newcommand{\R}{\mathbb R}
\newcommand{\Lip}{\operatorname{Lip}}
\newcommand{\diam}{\operatorname{diam}}
\newcommand{\de}{\delta}
\newcommand{\norm}[1]{\left\lVert #1\right\rVert}

\title{A Counterexample to Remark 3.3 in arXiv:1612.03808}
\author{fa\_banach\_001, agent\_lane\_12}
\date{2026-06-20}

\begin{document}
\maketitle

\section*{Status}

This packet gives a likely valid counterexample to the question in Remark 3.3
of Prochazka and Rueda Zoca, \emph{A characterisation of octahedrality in
Lipschitz-free spaces}, arXiv:1612.03808. The source paper asks whether
assertion (3) in that remark implies assertion (1). The answer is negative:
for \(r=2/5\), the metric space already appearing in Example 3.6 satisfies
assertion (3), in the only meaningful range \(0<\varepsilon<1-r\), but fails
assertion (1).

\section*{Source Passage}

Remark 3.3 states a chain of implications between four assertions, indexed by
\(0\le r<1\), and says that it is unknown whether assertion (3) implies
assertion (1).

\begin{center}
\includegraphics[width=0.96\textwidth]{figures/open_problem_crop.png}
\end{center}

For reference, the two assertions relevant here are:

\begin{description}
\item[(1)] For every \(\mu_1,\ldots,\mu_n\in S_{\F(M)}\) and every
\(\varepsilon>0\), there is \(u\ne v\in M\) such that
\[
 \left\|\mu_i+\frac{\de_u-\de_v}{d(u,v)}\right\|
 \ge 2-r-\varepsilon
\]
for all \(i\).

\item[(3)] For each \(\varepsilon>0\) and each finite \(N\subset M\), there
are \(u,v\in M\), \(u\ne v\), such that every 1-Lipschitz function
\(f:N\to\R\) admits an extension \(\widetilde f:M\to\R\) which is
\((1-r-\varepsilon)^{-1}\)-Lipschitz and satisfies
\[
 \widetilde f(u)-\widetilde f(v)\ge d(u,v).
\]
\end{description}

The denominator in (3) forces the usual small-\(\varepsilon\) reading
\(0<\varepsilon<1-r\). For larger \(\varepsilon\) the displayed Lipschitz
constant is not positive, so the original statement is not a meaningful
metric extension assertion.

\section*{Proof Intuition}

The relaxed extension assertion only needs two new points \(u,v\) on which
one can force a unit drop while remaining within a fairly large Lipschitz
constant. In the source paper's Example 3.6 metric, there are infinitely many
points \(x_n\) all at distance one from every old point except themselves, so
after a finite set \(N\) is chosen we can pick two fresh points \(x_i,x_j\).
Because the whole metric has diameter two, every 1-Lipschitz function on
\(N\) oscillates by at most two. For \(r=2/5\), the permitted constant is
larger than \(3/2\), exactly enough room to assign values \(a\) and \(a-1\)
to the two fresh points and then extend.

The free-space thickness assertion is stricter. Testing only the two opposite
molecules \(\mu=(\de_0-\de_z)/2\) and \(-\mu\) already blocks it: every
candidate molecule points too close to at least one of these two directions.
The Arens-Eells decompositions below show a uniform upper bound \(3/2\), while
assertion (1) would demand values arbitrarily close to \(2-r=8/5\).

\section*{The Counterexample}

Let
\[
 M=\{0,z\}\cup\{x_n:n\in\mathbb N\}
\]
with
\[
 d(0,z)=2,\qquad d(a,b)=1
\]
for every other pair of distinct points. This is a metric space; it is the
same metric used in Example 3.6 of the source paper.

\begin{theorem}
For \(r=2/5\), assertion (3) of Remark 3.3 holds for \(M\), but assertion
(1) fails for \(M\). Hence assertion (3) does not imply assertion (1).
\end{theorem}

\begin{proof}
First prove assertion (3). Let \(N\subset M\) be finite and let
\(0<\varepsilon<1-r=3/5\). Put
\[
 L=\frac1{1-r-\varepsilon}=\frac1{3/5-\varepsilon}.
\]
Then \(L>5/3>3/2\). Choose distinct \(u=x_i\) and \(v=x_j\) outside \(N\).
Thus \(d(u,v)=1\) and \(d(u,s)=d(v,s)=1\) for every \(s\in N\).

Let \(f:N\to\R\) be 1-Lipschitz. Since \(\diam M=2\), the oscillation
\(\operatorname{osc}_N f=\max_N f-\min_N f\) is at most \(2\). We need a real
number \(a\) such that
\[
 |a-f(s)|\le L,\qquad |a-1-f(s)|\le L\qquad(s\in N).
\]
The simultaneous conditions are equivalent to
\[
 a\in [\max_N f+1-L,\ \min_N f+L].
\]
This interval is nonempty because
\[
 \max_N f+1-L\le \min_N f+L
 \quad\Longleftrightarrow\quad
 1+\operatorname{osc}_N f\le 2L,
\]
and \(1+\operatorname{osc}_N f\le 3<2L\). Choose such an \(a\), and define
\[
 g|_N=f,\qquad g(u)=a,\qquad g(v)=a-1.
\]
The preceding inequalities control the pairs joining \(N\) to \(u\) or \(v\);
on \(N\) itself the map is \(L\)-Lipschitz because \(L>1\), and on the pair
\(\{u,v\}\) because \(|g(u)-g(v)|=1\le Ld(u,v)\). Hence \(g\) is
\(L\)-Lipschitz on \(N\cup\{u,v\}\), and \(g(u)-g(v)=1=d(u,v)\). The standard
formula
\[
 \widetilde g(p)=\inf_{q\in N\cup\{u,v\}}\bigl(g(q)+L\,d(p,q)\bigr)
\]
extends \(g\) to an \(L\)-Lipschitz function on all of \(M\). Therefore
assertion (3) holds.

Now prove that assertion (1) fails. Let
\[
 \mu=\frac{\de_0-\de_z}{d(0,z)}=\frac{\de_0-\de_z}{2}\in S_{\F(M)}.
\]
We will show that for every molecule
\[
 m_{uv}=\frac{\de_u-\de_v}{d(u,v)}
\]
one has
\[
 \min\{\norm{\mu+m_{uv}},\norm{-\mu+m_{uv}}\}\le \frac32.
\]
This uses only the elementary upper estimate for the Arens-Eells norm:
if \(\eta=\sum_k a_k(\de_{p_k}-\de_{q_k})\) with \(a_k\ge0\), then
\[
 \norm{\eta}\le \sum_k a_k\,d(p_k,q_k).
\]

If \(\{u,v\}=\{0,z\}\), then \(m_{uv}\) is \(\mu\) or \(-\mu\), so one of
\(\mu+m_{uv}\), \(-\mu+m_{uv}\) is zero. If exactly one endpoint is one of
\(\{0,z\}\) and the other is some \(x_n\), one of the two norms is at most
\(1\). For instance,
\[
 -\mu+m_{0x_n}
 =\frac12(\de_z-\de_{x_n})+\frac12(\de_0-\de_{x_n}),
\]
and both displayed distances are \(1\); the other orientations are identical
by symmetry.

It remains to consider \(u=x_i\), \(v=x_j\), \(i\ne j\). Then
\[
 \mu+m_{uv}
 =\frac12(\de_0-\de_v)+\frac12(\de_u-\de_z)
  +\frac12(\de_u-\de_v),
\]
and
\[
 -\mu+m_{uv}
 =\frac12(\de_z-\de_v)+\frac12(\de_u-\de_0)
  +\frac12(\de_u-\de_v).
\]
All distances appearing in these two decompositions are \(1\), so both norms
are at most \(3/2\). Hence the displayed minimum is always at most \(3/2\).

For \(r=2/5\), assertion (1) would require, for the pair \(\mu,-\mu\), that
there exist a molecule \(m_{uv}\) with
\[
 \norm{\mu+m_{uv}}\ge 2-r-\varepsilon,\qquad
 \norm{-\mu+m_{uv}}\ge 2-r-\varepsilon.
\]
Taking \(\varepsilon=1/20\), this lower bound is
\[
 2-\frac25-\frac1{20}=\frac{31}{20}>\frac32,
\]
contradicting the uniform estimate above. Thus assertion (1) fails.
\end{proof}

\section*{Computational Check}

The script \texttt{code/verify\_norm\_bounds.py} solves the finite
Arens-Eells transport linear programs for the four representative points
\(\{0,z,a,b\}\). It confirms
\[
 \max_{u\ne v}\min\{\norm{\mu+m_{uv}},\norm{-\mu+m_{uv}}\}=3/2.
\]
The proof above does not depend on this computation; it is included as an
independent sanity check of the molecule case analysis.

\section*{Novelty and Literature Check}

The run indexes \path{registry_index.tsv}, \path{solutions/index.tsv},
\path{attempts/index.tsv}, and \path{proof_gaps/index.tsv} were searched for
the arXiv id, title phrases, Remark 3.3, thick3, thick1, Whitley's thickness
index, and long trapezoid property phrases. No prior packet for this target
was found.

External web searches on 2026-06-20 for exact and near-exact phrases included:
\begin{itemize}
\item ``We do not know whether'' with ``thick3'' and ``thick1'';
\item ``assertion (3)'' and ``assertion (1)'' with ``Lipschitz-free'';
\item ``Whitley's thickness index'' with ``Lipschitz-free'' and ``long trapezoid''.
\end{itemize}
These searches returned the source paper or no answer to Remark 3.3. The
novelty confidence is therefore good but still subject to human literature
review.

\section*{References}

\begin{thebibliography}{9}
\bibitem{PRZ}
A. Prochazka and A. Rueda Zoca,
\emph{A characterisation of octahedrality in Lipschitz-free spaces},
arXiv:1612.03808, 2016.
\end{thebibliography}

\section*{Human Review Recommendation}

The recommended review focus is the small-\(\varepsilon\) convention in
assertion (3). Since the source's constant \((1-r-\varepsilon)^{-1}\) is
positive only for \(\varepsilon<1-r\), the counterexample is complete under
the necessary reading of the statement.

\end{document}
